\clearpage
\fancyhead[L]{\small INSTRUMENTE.}
\SourceLeaf{415}
\begin{center}
\Large INSTRUMENTE.
\end{center}

\begin{center}
Göttingische gelehrte Anzeigen. Stück 127. Seite 1257..1267. 1818. August 8.
\end{center}

\par\indent Von unserer \emph{neuen Sternwarte}, diesem grossen und sprechenden Denkmal der Liebe unserer Regierung für die Wissenschaft, ist bisher in unsern Blättern noch keine besondere Erwähnung geschehen, obgleich das Gebäude bereits seit anderthalb Jahren in dem Maasse vollendet ist, dass der Instrumentenvorrath der alten Sternwarte in dasselbe aufgenommen, und unterbrochen in einer Abtheilung des Gebäudes beobachtet werden konnte. Allein der Natur der Sache nach konnten diese Beobachtungen nur zu der Gattung derjenigen gehören, dergleichen auch auf der alten Sternwarte schon sich anstellen liessen, und es schien uns nicht passend, davon in diesen Blättern besondere Anzeige zu machen; verschiedene davon sind bereits in astronomischen Zeitschriften bekannt gemacht. Die neue Sternwarte, bestimmt, keiner nachzustehen, kann in den ihr gebührenden Rang erst durch den Besitz der \emph{festen Meridian-Instrumente} treten, und die Zeit ist jetzt nahe, wo sie vollständig ausgerüstet sein wird. --- Ueber das Aeussere des Gebäudes, welches der Würde seiner Bestimmung entspricht und der geschickten und geschmackvollen Ausführung unsers Hrn. Universitäts-Baumeisters \textsc{Müller} zur Ehre gereicht, werden wir hier nichts sagen. Auch eine vollständige Beschreibung der innern Einrichtung wird einem andern Orte vorbehalten bleiben. Aber die neuen Hauptinstrumente, welche dem Gebäude erst seinen wahren wissenschaftlichen Werth geben, sollen, so wie sie nach und nach ankommen und aufgestellt werden, durch diese Blätter näher angezeigt, und in die Bekanntschaft der Verehrer der Himmelskunde eingeführt werden. --- Das erste der neuen Meridianinstrumente, der \emph{Repsoldsche Meridiankreis}, kam im April d.\,J. an und wurde von dem Künstler selbst aufgestellt. Dieses bereits vor längerer Zeit von Hrn. \textsc{Repsold} in Hamburg ursprünglich zu seinem eignen Gebrauch verfertigte, und in dessen Privatsternwarte auf dem Hamburger Walle aufgestellt gewesene Instrument, ist den Freunden der Astronomie nicht unbekannt, indem verschiedene Auszüge aus Hrn. \textsc{Repsolds} Tagebüchern in der Monatlichen Correspondenz abgedruckt sind; auch in unsern Blättern (1811 S. 1290 [S. 327 d. B.]) haben wir einige von Hrn. Professor \textsc{Schumacher} an diesem Kreise gemachten Beobachtungen angezeigt. Da Hrn. \textsc{Repsolds} Sternwarte in der für Hamburg so unglücklichen Periode demolirt war, und der Besitzer demnach das Instrument unbenutzt lassen musste, so genehmigte unsere Regierung den Vorschlag des Hrn. Hofrath \textsc{Gauss}, dasselbe für unsere neue Sternwarte anzukaufen. Der Künstler übernahm dabei mehrere so wichtige und bedeutende Vervollkommnungen an dem Instrumente, dass dieses gewissermassen ein ganz neues geworden ist. Die durch manche unvorhergesehene Umstände oft unterbrochene Arbeit an diesen neuen Einrichtungen ist die Ursache, dass die Ablieferung des Instruments sich so lange verzögert hat, zu dessen Aufstellung auf der Sternwarte bereits seit anderthalb Jahren alles vorbereitet war. --- Die Einrichtung dieses Instruments verdient um so mehr eine umständlichere Beschreibung, da sie, ganz aus den eignen Ideen des genialen Künstlers hervorgegangen, bis jetzt einzig ist. Es vereinigt in sich ein vollkommnes Mittagsfernrohr mit einem Kreise, der dem neuen Greenwicher Mauerkreise sehr ähnlich ist. Das Fernrohr von $7\frac{1}{2}$ Pariser Fuss Länge, $6\frac{3}{4}$ Fuss Brennweite und 46 Linien Oeffnung, befindet sich an einer Axe von 4 Fuss Länge, die wie bei andern Mittagsfernröhren zwischen zwei steinernen Pfeilern von $6\frac{1}{2}$ Fuss Höhe und 22 Zoll im Quadrat, aufgehängt ist. Durch drei besondere Balancirungen wird sowohl die Biegung des Fernrohrs, als die der Axe aufgehoben und das Gewicht des Ganzen so weit getragen, dass die Pfannen nur einen Druck von einigen Lothen erleiden. Die Pfannen \SourceLeaf{416} selbst sind von Bergkrystall, die vollkommen cylindrischen und gleich dicken Zapfen von
Glockenmetall; an letztern kann, wenn sie nach langem Gebrauch durch die geringe Reibung doch etwas
abgenutzt werden sollten, durch eine besondere Einrichtung eine andere Stelle zum Aufliegen gebracht
werden. Die Ocularröhre ist parallel mit der Axe verschiebbar, um jeden Faden in die Mitte des
Gesichtsfeldes bringen zu können; die Beleuchtung der Fäden geht durch die Axe. Die stärkste
Vergrösserung, welche gewöhnlich gebraucht wird, ist eine 96malige. Bei günstiger Luft lassen sich
Sterne bis zur dritten Grösse bei Tage ohne Mühe beobachten. Die Hänglibelle, wodurch die Axe
horizontal gestellt wird, ist höchst sorgfältig gearbeitet, und lässt Theile von Secunden mit
Sicherheit erkennen. Den beiden Pfeilern dient Eine grosse Platte zur gemeinschaftlichen Unterlage,
welche selbst durch grosse Quadern 12 Fuss tief im Boden fundirt ist. Die Solidität dieser Einrichtung
hat sich seit der Aufstellung auf das vollkommenste bewährt; Azimuth, Horizontalität und
Collimationslinie zeigen bis jetzt wenigstens eine fast absolute Unwandelbarkeit. --- Auf der Axe,
nahe dem einen Ende, sitzt der Kreis fest, der bei der Theilung $3\frac{1}{2}$ Fuss im Durchmesser hält.
Die Theilung ist mit einer bewundernswürdigen Genauigkeit durch sehr saubere Striche von 5 zu 5
Minuten gemacht, während der Kreis auf der Axe selbst sass, wodurch alle Excentricität vermieden ist.
Zur Versendung musste er freilich wieder abgenommen werden; durch die an Ort und Stelle geschehene
Wiederaufsetzung ist nur eine fast unmerkliche Excentricität entstanden, die nach der sorgfältigen
Prüfung des Hrn. Hofr. \textsc{Gauss} $0''8$ beträgt. Die Ablesung geschieht durch drei vortreffliche
mikrometrische Mikroskope, deren Träger an den schiebbaren Theilen der Lager so befestigt sind, dass
ihre Entfernung von der Kreisfläche, auch wenn diese von der Ebene etwas abweichen sollte,
vollkommen beständig bleibt; das eine Mikroskop sitzt unten, die beiden andern seitwärts $90^\circ$
von jenem entfernt. Ein viertes festes Mikroskop anzubringen verstattete der Bau des Instruments nicht
(früher war nur ein einziges da). Da die Wirkung einer Excentricität durch die beiden Seitenmikroskope
ganz aufgehoben wird, so bleibt bei der Ablesung aller, nur der dritte Theil davon übrig, der nur auf
$0''3$ sich belaufen, und wo es nöthig scheint mit in Rechnung gebracht werden kann. Ein viertes
\emph{bewegliches} Mikroskop zur Prüfung der Theilung selbst, wird der Künstler auf den Wunsch des Hrn.
Hofr. \textsc{Gauss} noch nachliefern. Die Läufer in den Mikroskopen, welche durch äusserst
gleichförmige von allem todten Gange durch die bekannte \textsc{Ramsden}sche Erfindung frei gemachte
Schrauben bewegt werden, führen ein kreisrundes Loch, in welchem die Bissection der Theilstriche mit
einer solchen Schärfe erkannt werden kann, dass man auf Theile von Secunden sicher ist. Der Kreis hat
neun Speichen, zwischen einem Paar derselben befindet sich ein kupferner Cylinder, welcher durch eine
ebenfalls höchst vortreffliche Libelle auf eine halbe Secunde genau horizontal gestellt werden kann.
Hierdurch werden nicht allein die etwanigen Veränderungen in der Stellung der Mikroskope bemerkbar,
sondern, da die Neigung der Gesichtslinie gegen diesen Cylinder durch Beobachtungen in gewechselter
Lage des Instruments bekannt wird, so erhält man auch absolute Zenithdistanzen. --- Bei der Ankunft
des Instruments war das Fadennetz im Fernrohr noch dasselbe, welches bei der ersten Verfertigung
eingezogen war, und aus drei Verticalen, und zwei horizontalen $21''$ von einander abstehenden
Spinnenfäden bestand. Da Hr. \textsc{Repsold} Ende Mai wiederum nach Göttingen zurückkam, so
vertauschte er dieses Fadennetz mit einem neuen, wodurch das Pointiren an Genauigkeit beträchtlich
gewonnen hat. Das neue Fadennetz besteht aus fünf verticalen und zwei horizontalen ausgesuchten
Spinnenfäden; letztere sind vollkommen parallel, und stehen $12''7$ von einander ab. Beim Beobachten
wird in der Regel der Stern in die Mitte zwischen diese beiden Fäden gebracht, welches mit sehr
grosser Schärfe geschehen, und wobei man die Antritte an die verticalen Fäden immer leicht und scharf
beobachten kann. Die Zwischenzeiten zwischen den einzelnen Verticalfäden betragen für Sterne im
Aequator etwas über 16 Secunden. --- Die Beobachtungen des Polarsterns gehören auf jeder wohlbestellten \SourceLeaf{417} Sternwarte aus bekannten Gründen zu den täglichen; bei einem neuen Mittagsfernrohr sind sie
aber doppelt wichtig, da auf dieselben die Hauptberichtigungen gegründet werden müssen. Hr. Hofr.
\textsc{Gauss} hat daher, so oft es das im Juni sehr günstige Wetter erlaubte, alle obern und untern Culminationen
dieses Sterns sorgfältig beobachtet. Wir geben hier als eine Probe von dem, was das Instrument
leistet, die bisher beobachteten geraden Aufsteigungen, nebst der Vergleichung mit der von Hrn.
\textsc{Bessel} im astronomischen Jahrbuche 1817 gegebenen und im folgenden Jahrgange verbesserten Hülfstafel.
Der Bruch des Tages ist von der untern Culmination, welche zu dem angesetzten Datum gehört,
gezählt. Die letzte Columne drückt das Quadrat der Genauigkeit jedes einzelnen Resultats aus, nach
Massgabe der Anzahl der concurrirenden Fadenantritte, welches Quadrat von einigen Astronomen ganz
schicklich das Gewicht der Beobachtung genannt ist; als Einheit liegt dabei diejenige Genauigkeit zum
Grunde, die eine aus zwei einander folgenden Culminationen, die jede nur an Einem Faden beobachtet
sind, geschlossene Bestimmung hat. Die folgenden 19 Bestimmungen gründeten sich auf 41 beobachtete
Culminationen, von denen immer zwei oder drei unmittelbar auf einander folgende zu Einem Resultate
verbunden wurden.

\begin{center}
\begin{tabular}{r@{\hspace{4pt}}l|c c c|c|c}
\multicolumn{2}{c}{1818} & \multicolumn{3}{c}{\shortstack{Beob. gerade\\Aufsteigung}} & \shortstack{Correction\\der Tafel} & Gewicht \\
\hline
Juni & 3.25  & $0^{\mathrm{u}}$ & $55^{\mathrm{m}}$ & $58^{\mathrm{s}}65$ & $+3^{\mathrm{s}}93$ & 2.67 \\
     & 4.25  & & & 58.19 & $+2.80$ & 5.00 \\
     & 5.25  & & & 58.56 & $+2.45$ & 5.00 \\
     & 6.25  & & & 58.91 & $+2.22$ & 4.00 \\
     & 8.50  & $0$ & $56$ & $1.14$ & $+2.97$ & 6.40 \\
     & 10.25 & & & 2.58 & $+3.20$ & 3.75 \\
     & 11.25 & & & 2.68 & $+2.59$ & 1.67 \\
     & 12.50 & & & 3.31 & $+2.36$ & 4.29 \\
     & 16.75 & & & 6.11 & $+2.13$ & 4.44 \\
     & 18.25 & & & 7.31 & $+2.26$ & 5.00 \\
     & 19.25 & & & 8.29 & $+2.51$ & 2.86 \\
     & 21.25 & & & 10.49 & $+3.26$ & 5.00 \\
     & 22.75 & & & 12.02 & $+3.60$ & 1.33 \\
     & 26.25 & & & 13.49 & $+2\ 60$ & 1.60 \\ 
     & 27.25 & & & 14.77 & $+3.14$ & 4.44 \\
     & 29.50 & & & 16.71 & $+3.42$ & 3.43 \\
Juli & 7.50  & & & 22.51 & $+3.31$ & 6.00 \\
     & 8.25  & & & 22.97 & $+3.22$ & 1.67 \\
     & 9.25  & & & 23.89 & $+3.40$ & 3.75 \\
\end{tabular}
\end{center}

\par\indent Das Mittel der durch diese 19 Beobachtungen gefundenen Correctionen der Tafel ist $+2^{\mathrm{s}}88$. In
so fern man sich erlaubt, den Tafelfehler während dieser Zeit als beständig zu betrachten, gibt die Anwendung
der Wahrscheinlichkeitstheorie des Hrn. Hofrath \textsc{Gauss} den wahrscheinlichen Fehler einer
Beobachtung, deren Genauigkeit 1 ist, $=0^{\mathrm{s}}64$, und den wahrscheinlichen Fehler des Endresultats
$=0^{\mathrm{s}}076$. Wir fügen ferner, als eine zweite Probe, die Beobachtungen des Uranus bei, welche Hr. Hofr.
\textsc{Gauss} zur Zeit der Opposition dieses Planeten angestellt hat. Die Declinationen gründen sich auf die
Vergleichung mit der unmittelbar vorausgegangenen untern und der unmittelbar folgenden obern Culmination
des Nordsterns, blos die vom 7ten und 13ten Juni ausgenommen, wo blos die untere Culmination
beobachtet war, und wo die Declination des Nordsterns aus \textsc{Bessels} Tafel zum Grunde gelegt
ist. Die Refractionen sind aus \textsc{Bessels} Tafel, die sich auf die \textsc{Bradley}schen Beobachtungen gründet,
genommen; würden dieselben aus \textsc{Carlinis} Tafel entlehnt, so würden die Declinationen des Uranus $3''$
südlicher ausfallen. Wenn in Zukunft erst eine hinlängliche Anzahl schicklicher Beobachtungen beisammen
sein wird, wird sich zeigen, welche von beiden Tafeln dem hiesigen Clima am angemessensten ist.

\SourceLeaf{418}
\begin{center}
\setlength{\tabcolsep}{3pt}
\begin{tabular}{l|ccc|ccc|ccc}
\hline
1818 & \multicolumn{3}{c|}{Mittlere Zeit} & \multicolumn{3}{c|}{Gerade Aufsteig.} & \multicolumn{3}{c}{Südl. Abweich.} \\
\hline
Juni 3 & $12^{\mathrm u}$ & $22^{\mathrm m}$ & $17^{\mathrm s}5$ & $257^\circ$ & $23'$ & $28''3$ & $23^\circ$ & $1'$ & $41''5$ \\
4 & 12 & 18 & 10.9 & & 20 & 47.8 & & 1 & 29.0 \\
5 & 12 & 14 & 4.6 & & 18 & 9.9 & & 1 & 14.4 \\
7 & 12 & 5 & 51.6 & & 12 & 54.4 & & 0 & 50.5 \\
8 & 12 & 1 & 45.1 & & 10 & 15.0 & & 0 & 38.4 \\
9 & 11 & 57 & 38.6 & & 7 & 36.4 & & 0 & 27.4 \\
10 & 11 & 53 & 32.1 & & 4 & 56.5 & & 0 & 12.6 \\
11 & 11 & 49 & 25.7 & & 2 & 18.0 & & 0 & 2.9 \\
12 & 11 & 45 & 19.2 & 256 & 59 & 33.9 & 22 & 59 & $50\;9$ \\ 
13 & 11 & 41 & 12.7 & & 57 & 0.4 & & 59 & 40.5 \\
\hline
\end{tabular}
\end{center}

\Needspace{18\baselineskip}\par\indent Hr. Hofr. \textsc{Gauss} übertrug die Vergleichung dieser Beobachtungen mit den \textsc{Delambre}schen Uranustafeln Hrn. \textsc{Dirksen}, welcher sich bei uns dem Studium der Astronomie mit ausgezeichnetem Eifer widmet. Die Resultate dieser Vergleichung sind folgende:

\begin{center}
\begin{tabular}{c|cc}
& \multicolumn{2}{c}{Unterschied} \\
& Ger. Aufst. & Abweichung \\
\hline
Juni 3 & $-54.1$ & $+10.3$ \\
4 & $-52.1$ & $+10.6$ \\
5 & $-52.9$ & $+12.9$ \\
7 & $-55.8$ & $+12.3$ \\
8 & $-55.9$ & $+12.0$ \\
9 & $-56\;8$ & $+10.5$ \\
10 & $-56.4$ & $+12.9$ \\
11 & $-57.2$ & $+10.1$ \\
12 & $-52.4$ & $+9\;4$ \\ 
13 & $-58.1$ & $+7.4$ \\
\hline
Mittel & $-55.2$ & $+10.8$ \\
\end{tabular}
\end{center}

\Needspace{15\baselineskip}\par\indent Hieraus ergab sich, nach Hrn. \textsc{Dirksens} Rechnung:

\begin{center}
Zeit der Opposition 1818 Juni 9.\\
$5^{\mathrm u}\,30^{\mathrm m}\,43^{\mathrm s}$ Mittlere Zeit in Göttingen
\end{center}
\begin{center}
\begin{tabular}{ll}
Wahre Länge . . . . . . . . . . . . . . & $258^\circ\;10'\;34''5$ \\
Geocentrische Breite . . . . . . . . . . . . & $0\;4\;12.1$ \\
Heliocentrische . . . . . . . . . . . . . . & $0\;3\;58.7$ Südl. \\
\end{tabular}
\end{center}

\noindent und der Fehler der \textsc{Delambre}schen Tafeln

\begin{center}
\begin{tabular}{ll}
in der Länge . . . . . . . . . . . . . . & $-45''3$ \\
in Heliocentrischer Breite . . . . . . . . . . & $+14.6$ \\
\end{tabular}
\end{center}

\par\indent Wir können diese Anzeige nicht schliessen, ohne zugleich der \textsc{Shelton}schen Pendeluhr zu erwähnen, die zwar schon seit 47 Jahren auf der alten Sternwarte im Gebrauch gewesen war, deren Vortrefflichkeit aber erst jetzt ganz gewürdigt werden kann, da früher die Sternwarte kein Mittel zur allerschärfsten Zeitbestimmung besass. Wir fügen zum Beweise nur das Register ihres Ganges von einem Monat bei, und bemerken, dass dieselbe Gleichförmigkeit während der ganzen Zeit statt gefunden hat, wo das \textsc{Repsold}sche Instrument zur Bestimmung des Ganges gedient hat. Das Datum bezieht sich auf Sonnentage, und die Stunden der Uhrzeit sind, um die Uebersicht zu erleichtern jedesmal über 24 hinaus bis zum nächsten Mittage fortgezählt; die letzte Columne enthält die Anzahl der Sterne des \textsc{Maskelyne}schen Fundamentalcatalogs, auf welche die Zeitbestimmung gegründet ist.

\SourceLeaf{419}
\begin{center}
\begin{tabular}{ll|c|c|c}
 & & Uhrzeit & \shortstack{Stand gegen\\Sternzeit} & \shortstack{Verglichene\\Sterne} \\
\hline
Juni & 15 & $15^{\mathrm u}\;22^{\mathrm m}$ & $+0^{\mathrm s}41$ & 4 \\
 & 17 & $20\;\;50$ & $+0.28$ & 2 \\
 & 18 & $11\;\;37$ & $+0.29$ & 2 \\
 & 20 & $12\;\;28$ & $+0.11$ & 2 \\
 & 21 & $11\;\;37$ & $-0.06$ & 2 \\
 & 22 & $7\;\;\;30$ & $-0.59$ & 1 \\
 & 23 & $11\;\;40$ & $-0.90$ & 1 \\
 & 26 & $14\;\;54$ & $-1.78$ & 5 \\
 & 27 & $16\;\;15$ & $-1.86$ & 4 \\
 & 28 & $14\;\;41$ & $-1.99$ & 1 \\
 & 29 & $18\;\;\;\;1$ & $-2.32$ & 6 \\
 & 30 & $16\;\;\;\;9$ & $-2.43$ & 5 \\
\end{tabular}
\end{center}

\begin{center}
\begin{tabular}{ll|c|c|c}
 & & Uhrzeit & \shortstack{Stand gegen\\Sternzeit} & \shortstack{Verglichene\\Sterne} \\
\hline
Juli & 1 & $14^{\mathrm u}\;7^{\mathrm m}$ & $-2^{\mathrm s}92$ & 1 \\
 & 2 & $14\;\;51$ & $-3.11$ & 2 \\
 & 6 & $25\;\;17$ & $-3.87$ & 4 \\
 & 7 & $22\;\;31$ & $-3.90$ & 7 \\
 & 8 & $16\;\;36$ & $-4.03$ & 4 \\
 & 9 & $28\;\;\;7$ & $-4.21$ & 4 \\
 & 10 & $15\;\;\;3$ & $-4.28$ & 3 \\
 & 13 & $29\;\;24$ & $-4.53$ & 2 \\
 & 14 & $29\;\;14$ & $-4.62$ & 4 \\
 & 15 & $15\;\;56$ & $-4.73$ & 5 \\
 & 16 & $14\;\;41$ & $-4.93$ & 1 \\
\end{tabular}
\end{center}

\par\indent Hiebei darf nicht unerwähnt bleiben, dass die Aufstellung der Uhr erst noch provisorisch ist. Sie ruht zwar auf einem besondern steinernen Fundament, aber nur vermittelst ihres Gehäuses, ohne unmittelbar an den steinernen Pfeiler, neben welchem dieses gestellt ist, befestigt zu sein. Sie geht einen Monat in Einem Aufzuge; um die Mitte dieser Zeit kommt das Gewicht der Linse gegenüber zu stehen, welches man sonst wohl für nachtheilig gehalten und durch mancherlei Einrichtungen zu vermeiden gesucht hat. Allein das Register der \textsc{Shelton}schen Uhr zeigt um diese Zeit gar keine spürbare Veränderung im Gange, der überhaupt viel regelmässiger ist, als der mancher andern Uhren auf den ersten Sternwarten, auch solcher, die man unmittelbar an einen steinernen Pfeiler befestigt hat. Sollte man nicht hieraus schliessen, dass beide Umstände von viel geringerer Wichtigkeit sind, als man bisher gewöhnlich geglaubt hat?
